\title{SMC's Tidal Disruption as a New Dark Matter Probe with Spec-S5}

\author{Himansh Rathore}
\affil{Steward Observatory, University of Arizona}

\author{Gurtina Besla}
\affil{Steward Observatory, University of Arizona}

\author{Hayden R. Foote}
\affil{Steward Observatory, University of Arizona}

\textbf{Facts:} The LMC-SMC system is the nearest example of an interacting galaxy pair. In particular, the SMC has directly collided (impact parameter $\sim 2$~kpc) with the LMC $<200$~Myr ago. As such, the SMC's stellar disk is disrupting due to the LMC's gravitational tides (REFS). 

\textbf{Importance of Facts:} The process of tidal disruption is sensitive to the dark matter profile of the disrupting galaxy [REFS], which in turn is sensitive to the interactions between dark matter particles [REFS]. As such, the SMC's tidal disruption is a new probe of dark matter microphysics.

\textbf{Problem:} Precise measurements of the velocity field of the SMC's old stellar populations, like red giant stars (G $\sim 17$ mag) and red clump stars (G $\sim 20$ mag) are required to constrain the SMC's tidal disruption radius. The measurement precision needs to be high enough that such comparisons between different dark matter models is possible (REFS). In particular, the SMC's azimuthally averaged velocity profile needs to be measured to a precision of at least 1 km/s to achieve a measurement uncertainty in the tidal radius that is comparable to the typical spatial resolution of simulations (< 0.1 kpc, [REFS]). Although Gaia has made significant strides in measuring the SMC's red giant proper motion field, the azimuthally averaged velocity field still has an uncertainty of $\sim 6-10$ km/s [REF], which translates to an uncertainty in the tidal radius of $>1$~kpc [REFS] --- at least an order of magnitude worse than the required precision. Current line of sight (LoS) velocity measurements (from e.g. BOSS, APOGEE-2) are also not sufficient to precisely measure the SMC's tidal radius [Zivick+2021].

\textbf{Solution:}

\textbf{Importance of Solution:}

\textbf{Broader Impacts}:  
